\documentclass[a4paper,14pt]{extarticle}

\usepackage{geometry}
\geometry{
  a4paper,
  left=2.5cm,
  right=2.5cm,
  top=2.5cm,
  bottom=2.5cm,
  bindingoffset=0cm
}

\usepackage{setspace}
\onehalfspacing

% Документ рассчитан на сборку XeLaTeX.
\usepackage{fontspec}
\IfFontExistsTF{Times New Roman}
  {\setmainfont{Times New Roman}}
  {\setmainfont{TeX Gyre Termes}}

\usepackage[russian]{babel}
\usepackage{indentfirst}
\setlength{\parindent}{1.25cm}
\setlength{\parskip}{0pt}
\setlength{\emergencystretch}{2em}

\usepackage{titlesec}
\titleformat{\section}{\normalfont\bfseries\centering}{\thesection}{1em}{}
\titleformat{\subsection}{\normalfont\bfseries\raggedright}{\thesubsection}{1.2em}{}
\titleformat{\subsubsection}{\normalfont\itshape\raggedright}{\thesubsubsection}{1em}{}
\titleformat{\paragraph}[block]{\normalfont\bfseries\raggedright}{}{0pt}{}
\titlespacing*{\paragraph}{0pt}{0.8\baselineskip}{0.35\baselineskip}

\usepackage{tocloft}
\renewcommand{\cftsecleader}{\cftdotfill{\cftdotsep}}

\usepackage{amsmath}
\usepackage{amssymb}
\usepackage{array}
\usepackage{booktabs}
\usepackage{longtable}
\usepackage{tabularx}
\usepackage{enumitem}
\setlist{nosep}

\usepackage{caption}
\captionsetup{justification=centering,labelsep=endash}

\usepackage{xcolor}
\usepackage{graphicx}
\usepackage{tikz}
\usetikzlibrary{arrows.meta,positioning,shapes.geometric,fit}
\usepackage{hyperref}
\hypersetup{
  colorlinks=true,
  linkcolor=black,
  citecolor=black,
  urlcolor=black,
  pdftitle={Отчёт по учебной практике},
  pdfauthor={Кондрашов Даниил Владиславович},
  pdfborder={0 0 0}
}

\newcommand{\Hnull}{\ensuremath{H_0}}
\newcommand{\Hone}{\ensuremath{H_1}}
\newcommand{\DeltaBIC}{\ensuremath{\Delta\mathrm{BIC}}}
\newcommand{\code}[1]{\texttt{#1}}

\begin{document}

% Титульный лист. Формулировку вида практики при необходимости можно
% заменить без изменения остального документа.
\begin{titlepage}
  \begin{center}
    \vspace*{1cm}

    {\large \textbf{Санкт-Петербургский государственный университет}}

    \vspace{0.5cm}

    {\large Искусственный интеллект и наука о данных}

    \vspace{4.5cm}

    {\large Кондрашов Даниил Владиславович}

    \vspace{1cm}

    {\Large \textbf{Исследование пересмотра физических гипотез большими языковыми моделями под влиянием экспериментальных данных}}

    \vspace{1cm}

    {\Large \textbf{Отчёт по учебной практике}}

    \vspace{1.8cm}

    \begin{flushright}
      {Научный руководитель:} \\
      {Азимов Рустам Шухратуллович}
    \end{flushright}

    \vfill

    {\large Санкт-Петербург \\ 2026}
  \end{center}
\end{titlepage}

\pagenumbering{gobble}
\tableofcontents
\newpage
\pagenumbering{arabic}

\section{Введение}

Большие языковые модели (Large Language Models, LLM) всё чаще рассматриваются не только как средства генерации текста, но и как компоненты систем научного поиска. Такие системы могут формулировать гипотезы, предлагать эксперименты, интерпретировать наблюдения и строить объяснения. Современные бенчмарки уже моделируют полный исследовательский цикл в интерактивных физических средах, однако показывают, что даже сильные модели нестабильно обнаруживают скрытую структуру мира, а высокая точность предсказания не всегда сопровождается корректным объяснением закона~\cite{discoverphysics,physgym}. Поэтому для использования LLM в исследовательских задачах важно оценивать не только качество окончательного ответа, но и способность модели изменять выводы при появлении новых данных.

Пересмотр гипотезы является принципиальной частью научного рассуждения. Исследователь должен сохранять рабочую гипотезу, пока данные с ней согласуются, и отказываться от неё, когда накопленное свидетельство становится достаточно сильным. Для языковой модели эта задача осложняется влиянием ранее сформулированного ответа, заданного контекста и уверенности. Исследования belief revision показывают, что модели нередко ошибаются как при необходимом обновлении вывода, так и при его обоснованном сохранении~\cite{beliefrevision}. Другие работы выявляют одновременно действующие смещения: стремление поддерживать уже показанный собственный ответ и повышенную чувствительность к противоречащей внешней рекомендации~\cite{competingbiases}.

Существующие исследования обычно представляют новое свидетельство в виде естественно-языковых посылок, советов или фрагментов научных статей. При этом необходимость пересмотра часто задаётся категориально, а не через численную характеристику того, насколько наблюдения различают конкурирующие гипотезы. Из-за этого трудно отделить три возможных источника ошибки: неспособность извлечь статистическое свидетельство из числовых данных, некорректное объединение свидетельства с исходным убеждением и влияние формата запроса. Исследования байесовских свойств LLM подтверждают, что рост выраженной уверенности не всегда означает рост точности, а реакция на релевантность и надёжность контекста зависит от модели и набора задач~\cite{evidencetobelief}.

В данной работе рассматривается контролируемая физическая постановка с двумя вложенными гипотезами. Для каждого набора синтетических наблюдений независимо от ответа LLM вычисляется статистическая мера поддержки гипотез, а режимы генерации калибруются методом Монте-Карло. Затем один и тот же набор данных предъявляется модели при нейтральном и заведомо ошибочном исходном распределении вероятностей гипотез. Такой дизайн позволяет исследовать, растёт ли вероятность правильного выбора вместе с объективной доказательной силой и насколько дополнительное свидетельство требуется модели для преодоления ошибочного исходного убеждения.

На текущем этапе выполнены анализ предметной области, постановка исследовательских вопросов, разработка протокола первого эксперимента, программная реализация воспроизводимого конвейера и полный запуск двух локальных LLM. Обе модели завершили по 2400 испытаний без ошибок разбора, но во всех условиях выбирали \Hone. Поэтому полученная точность 50\% отражает не случайное угадывание, а вырожденную стратегию постоянного выбора более сложной гипотезы. Дополнительный анализ показывает, что знак \DeltaBIC{} указывает истинную гипотезу в 86\% уникальных наборов, однако модели не используют этот сигнал и не меняют ответ при изменении исходного распределения вероятностей.

\section{Цели и задачи работы}

\subsection{Цель работы}

Цель магистерской работы состоит в исследовании способности больших языковых моделей пересматривать предварительно заданную физическую гипотезу при предъявлении экспериментальных данных различной объективной доказательной силы, а также в выявлении этапа, на котором преимущественно возникает расхождение с нормативным статистическим решением.

\paragraph{Основной исследовательский вопрос}

\begin{quote}
Насколько вероятность пересмотра ошибочной физической гипотезы языковой моделью соответствует объективной доказательной силе экспериментальных данных?
\end{quote}

\paragraph{Дополнительный исследовательский вопрос}

\begin{quote}
Возникает ли ошибка преимущественно при извлечении статистического свидетельства из числовых данных или при интеграции уже явно предъявленного свидетельства с исходным убеждением?
\end{quote}

\subsection{Цель практики}

Целью практики является проведение первого этапа исследования: от анализа ближайших работ и операционализации понятий до реализации воспроизводимого конвейера, запуска двух локальных LLM и анализа результатов первого контролируемого эксперимента.

\subsection{Задачи практики}

Для достижения поставленной цели необходимо решить следующие задачи:

\begin{enumerate}
  \item проанализировать исследования пересмотра убеждений LLM, влияния исходной уверенности и контекста, а также применения языковых моделей к научному поиску;
  \item определить исследовательский пробел и сформулировать основной и дополнительный исследовательские вопросы;
  \item разработать физическую постановку с известной истинной гипотезой и независимо измеряемой доказательной силой наблюдений;
  \item определить экспериментальные факторы, контрольные условия, объём выборки и метрики оценки;
  \item реализовать генерацию и калибровку данных, формирование запросов, обращение к локальным LLM, разбор ответов и сохранение артефактов запуска;
  \item обеспечить воспроизводимость исследования за счёт версионирования конфигураций, запросов, наборов данных и сведений о моделях;
  \item провести основные запуски двух локальных LLM, проанализировать результаты и определить необходимые изменения постановки перед проверкой разных форм предъявления свидетельства.
\end{enumerate}

\section{Обзор предметной области}

\subsection{Границы области и состав обзора}

Обзор посвящён корректности формирования и пересмотра гипотез языковой моделью после получения новых данных. Термин «убеждение» используется здесь операционально: это выбранный моделью ответ и выраженная ею уверенность, а не устойчивое внутреннее состояние, память между запросами или знание, содержащееся в весах. Такое ограничение необходимо, поскольку по наблюдаемым ответам нельзя непосредственно восстановить внутренний механизм модели.

В рамках работы подробно рассмотрены ровно восемь исследований:

\begin{sloppypar}
\begin{enumerate}
  \item \textit{Anchoring Depends on Confidence and Post-Training in Language Models};
  \item \textit{Belief Revision: The Adaptability of Large Language Models Reasoning};
  \item \textit{Competing Biases underlie Overconfidence and Underconfidence in LLMs};
  \item \textit{DiscoverPhysics: Benchmarking LLMs for Out-of-the-Box Scientific Thinking};
  \item \textit{Evolving Scientific Discovery by Unifying Data and Background Knowledge with AI Hilbert};
  \item \textit{Experiments or Outcomes? Probing Scientific Feasibility in Large Language Models};
  \item \textit{From Evidence to Belief: A Bayesian Epistemology Approach to Language Models};
  \item \textit{PhysGym: Benchmarking LLMs in Interactive Physics Discovery with Controlled Priors}.
\end{enumerate}
\end{sloppypar}

Пять работ исследуют более узкие эффекты обновления ответа, уверенности и контекста; \textsc{DiscoverPhysics} и \textsc{PhysGym} рассматривают LLM в полном цикле физического исследования; AI Hilbert не использует LLM, но даёт важную формальную схему объединения исходной теории и наблюдений. Источники по BIC, доверительным интервалам и инфраструктуре vLLM, приведённые далее в списке литературы, являются вспомогательными техническими ссылками и не входят в эти восемь работ обзора.

\subsection{Пересмотр вывода и сохранение убеждения}

В работе Belief Revision предложен набор Belief-R для проверки того, умеет ли модель изменить исходный вывод после появления значимой посылки и сохранить вывод, когда пересмотр не требуется~\cite{beliefrevision}. Сначала модель решает задачу на основе двух естественно-языковых посылок по схеме modus ponens или modus tollens. Затем добавляется третья посылка. В условии Belief Update прежний вывод требуется отозвать, а в условии Belief Maintain -- сохранить. Авторы отдельно проверяют качество исходного рассуждения и для основного анализа отбирают модели, достигшие не менее 75\% точности до обновления.

Эксперименты примерно с 30 моделями показывают, что даже сильные системы нестабильно определяют момент необходимого пересмотра. Значение агрегированной метрики BREU у крупных моделей обычно остаётся около 50\%, а улучшение обновления нередко сопровождается ухудшением сохранения. Modus tollens оказывается сложнее modus ponens; Chain-of-Thought и Plan-and-Solve не дают устойчивого улучшения для всех моделей. Эти результаты важны методически: высокая частота смены ответа сама по себе не является успехом. Необходимо проверять как обоснованный отказ от гипотезы, так и её обоснованное сохранение.

Вместе с тем Belief-R использует синтетические бытовые причинные посылки. Необходимость пересмотра определяется категориально на основе человеческой разметки, доказательная сила новой информации не задаётся численно, а уверенность модели не измеряется. Поэтому работа не отвечает на вопрос, насколько сильными должны быть физические данные, чтобы отказ от исходной гипотезы был обоснован.

Наиболее близкую количественную схему предлагает работа Competing Biases~\cite{competingbiases}. В её двухэтапной процедуре LLM сначала отвечает на вопрос и выражает исходную уверенность. Затем в независимом запросе первоначальный ответ показывается или скрывается, а условный советчик поддерживает либо опровергает его. Надёжность советчика задаётся численно, преимущественно в диапазоне от 50\% до 100\%, а наблюдаемое обновление сравнивается с Bayesian ideal observer.

Авторы обнаруживают два одновременно действующих эффекта. Повторный показ собственного ответа повышает уверенность в нём и снижает частоту смены выбора даже без новой содержательной информации. Одновременно противоречащая рекомендация получает избыточный вес относительно байесовского ориентира. Следовательно, одно и то же итоговое решение может быть результатом конкуренции между поддержанием собственного выбора и повышенной чувствительностью к противоречию. Для настоящей работы отсюда следуют необходимость независимых контекстов, нейтрального условия и раздельного учёта исходного убеждения и направления свидетельства.

Ограничение этой постановки состоит в природе свидетельства: реального советчика нет, его ответ и заявленную точность задаёт экспериментатор. Исходная уверенность LLM измеряется, но не изменяется как независимый фактор. Модель не извлекает статистическое свидетельство из наблюдений, поэтому остаётся открытым вопрос, как эти эффекты проявятся при работе с физическими числовыми данными.

\subsection{Уверенность, числовые якоря и надёжность контекста}

Работа Anchoring Depends рассматривает влияние нерелевантного числа на количественный ответ LLM~\cite{anchoring}. Авторы проверяют шесть базовых, instruction-tuned и дистиллированных моделей семейств Llama и Qwen. В основном наборе из 50 вопросов каждый пример предъявляется без якоря, с низким и с высоким якорем. Распределения строятся по девяти фиксированным кандидатам, а величина сдвига измеряется расстоянием по полной вариации.

Чем выше исходная определённость распределения модели, тем слабее воздействие якоря. После учёта определённости фактическая правильность исходного ответа не даёт существенного дополнительного объяснения: уверенная ошибка может быть столь же устойчивой, как верный ответ. Высокие якоря в среднем действуют сильнее низких, а характер зависимости меняется после дообучения. Однако уверенность в этой работе наблюдается, а не задаётся экспериментатором, и используется фиксированный набор ответов. Нерелевантный якорь также принципиально отличается от содержательных данных, которые могут обоснованно требовать пересмотра. Для настоящего исследования работа мотивирует числовой контроль, позволяющий отличить реакцию на доказательную силу от реакции на само присутствие чисел.

From Evidence to Belief сопоставляет поведение LLM с принципами байесовского подтверждения, опровержения и нерелевантности~\cite{evidencetobelief}. На задачах SciQ, TriviaQA и GSM8K сравниваются отсутствие контекста, эталонные сведения, ложная, неполная, противоречивая и нерелевантная информация. В отдельной серии варьируются авторитетность источника, подробность, свежесть и описание способа получения данных. Уверенность измеряется словесной оценкой, вероятностями токенов и повторным сэмплированием, а отдельно рассчитываются точность и ошибка калибровки.

Эталонный контекст обычно повышает точность и уверенность, но не гарантирует улучшения калибровки. Ложные и нерелевантные сведения могут ухудшать качество, причём внешне близкий, но нерелевантный текст иногда мешает сильнее очевидно постороннего. Более убедительная подача способна повысить уверенность без сопоставимого роста точности. Эти результаты требуют раздельно измерять выбор, уверенность и калибровку. Однако «сила» свидетельства в работе задаётся качественными свойствами текста: отношение правдоподобий и ожидаемая величина обновления для конкретного примера не рассчитываются.

Обе работы показывают, что уверенность нельзя автоматически считать показателем правильности. Anchoring Depends отделяет уверенность от точности, а From Evidence to Belief -- уверенность от калибровки и фактического качества ответа. Для проектируемого эксперимента это означает, что итоговая гипотеза и выраженная уверенность должны анализироваться как разные зависимые переменные.

\subsection{Экспериментальный контекст и научная допустимость}

В работе Experiments or Outcomes? модель оценивает научную гипотезу как допустимую или недопустимую и формирует краткое обоснование~\cite{experimentsoutcomes}. Сравниваются четыре условия: только гипотеза, гипотеза с описаниями экспериментов, с их результатами и с обоими типами контекста. Полнота информации принимает значения 0, 0{,}5 и 1, а частичные наборы формируются случайно и усредняются по пяти запускам.

На наборе MATTER-OF-FACT результаты экспериментов часто оказываются полезнее их описаний, но эффект зависит от модели. Например, для Gemini Pro точность составляет 0{,}67 без контекста, 0{,}53 с полным описанием экспериментов и 0{,}66 с полными результатами; для GPT-5.1 описания в одной из постановок, напротив, оказываются полезнее результатов. Частичный контекст способен ухудшить решение сильнее полного, тогда как на наборе REASONS даже неполные сведения нередко помогают. Следовательно, дополнительный контекст не является автоматически полезным, а совместное предъявление описаний и результатов не гарантирует преимущества.

Эта работа близка по различению формы свидетельства, но описания экспериментов и их результаты содержат разную информацию. Их сравнение не эквивалентно сравнению сырых измерений и статистической обработки одной выборки. Кроме того, полнота контекста не является объективной доказательной силой, а остаточный риск утечки знаний и использование ROUGE для обоснований ограничивают интерпретацию. В настоящей работе разные формы свидетельства должны строиться из одних и тех же наблюдений.

\subsection{LLM в интерактивном физическом исследовании}

\textsc{DiscoverPhysics} оценивает полный цикл исследования неизвестной физической системы~\cite{discoverphysics}. Бенчмарк содержит 22 симулированных мира с законами, намеренно отличающимися от привычной физики. Агент выбирает условия эксперимента, получает зашумлённые траектории, уточняет гипотезу, а затем представляет Python-реализацию закона и текстовое объяснение. Одиннадцать моделей проверяются по пять раз на каждом мире; оценка включает ошибку траектории, качество объяснения и pass@k.

Работа показывает расхождение между предсказательной точностью и концептуальным пониманием. GPT-5.5 быстрее снижает MSE, но после нескольких раундов улучшение подгонки слабо сопровождается улучшением объяснения. Claude Opus 4.7 лучше использует дополнительные раунды и достигает более высокого качества объяснения и pass@k, хотя его ошибка предсказания выше. Даже лучший pass@5 составляет около 50\% миров, то есть успех хотя бы в одной из пяти попыток, а не вероятность успеха одного запуска. Особенно трудны миры со скрытыми видами частиц и другой латентной структурой.

\textsc{PhysGym} содержит 97 задач из шести областей физики и управляет объёмом предварительной информации~\cite{physgym}. На четырёх уровнях последовательно скрываются контекст задачи, физический смысл и привычные обозначения переменных. Агент выбирает эксперименты, получает значения симулятора и сохраняет историю наблюдений и гипотез. Сокращение подсказок в среднем ухудшает качество: например, доля успешных решений Gemini Pro снижается с 74{,}23\% на первом уровне до 31{,}96\% на четвёртом. Однако зависимость не всегда монотонна, а среди успешных запусков при меньшем объёме контекста модели нередко проводят больше экспериментов и предлагают больше гипотез.

Оба бенчмарка полезны тем, что используют среды с известным истинным законом и сохраняют процесс исследования. Вместе с тем они одновременно проверяют выбор эксперимента, работу с шумом, поиск класса закона, уточнение параметров и формирование объяснения. Ошибку нельзя однозначно связать с нежеланием пересматривать гипотезу. Доказательная сила каждого набора наблюдений относительно двух фиксированных альтернатив также не рассчитывается. Поэтому для изоляции пересмотра в настоящей работе всем условиям предъявляются одинаковые наблюдения, а выбор эксперимента исключён из первого дизайна.

\subsection{Формальное объединение теории и данных}

AI Hilbert решает задачу научного открытия средствами полиномиальной оптимизации, а не языковой моделью~\cite{aihilbert}. На вход подаются числовые наблюдения, аксиомы в виде равенств и неравенств и ограничения на сложность формулы. Система ищет закон, минимизируя ошибку на данных и отклонение от следствий исходной теории, а при возможности строит алгебраический сертификат выводимости.

Эксперименты показывают, что теория и данные способны взаимно компенсировать недостаток друг друга. При полной теории некоторые законы выводятся без наблюдений; при противоречивых аксиомах данные помогают исключить неверные предпосылки. В синтетическом опыте с третьим законом Кеплера после удаления одной аксиомы потребовалось 46 точек, а после удаления всех -- 71. Эти значения относятся к конкретной постановке без шума и не являются универсальной количественной ценой исходного знания.

Для настоящей работы особенно важны два вывода. Во-первых, исходную теорию, наблюдения и процедуру выбора закона следует разделять явно. Во-вторых, малая ошибка на имеющихся точках не доказывает восстановление правильной зависимости. Однако вес соответствия теории в целевой функции AI Hilbert не является prior в байесовском смысле, а расстояние до теории не является доказательной силой данных. Работа не исследует вероятность пересмотра убеждения LLM.

\subsection{Выводы обзора и исследовательский пробел}

Восемь рассмотренных работ дают согласованную, но неполную картину. Модели могут ошибаться как при необходимом пересмотре, так и при сохранении вывода; повторный показ собственного ответа влияет на дальнейший выбор; уверенность, правильность и калибровка способны расходиться; нерелевантные числа и форма контекста меняют распределение ответов; хорошая подгонка к наблюдениям не гарантирует правильного закона. Одновременно сила и направление этих эффектов зависят от модели, задачи и экспериментального режима.

При этом широкие бенчмарки научного поиска не изолируют пересмотр от выбора эксперимента и поиска закона. Более узкие работы тщательно контролируют контекст, но используют категориальные посылки, искусственного советчика или качественные признаки надёжности. Они не позволяют определить, какой объективной силы должны достичь физические наблюдения при заданном исходном убеждении, чтобы пересмотр стал нормативно обоснованным. Кроме того, они не разделяют неспособность извлечь свидетельство из чисел и неспособность обновить вывод по уже рассчитанному свидетельству.

Научная специфика предлагаемой постановки состоит не в первом исследовании belief revision или LLM в физике, а в объединении следующих элементов:

\begin{enumerate}
  \item двух заранее фиксированных физических гипотез, одна из которых действительно порождает данные;
  \item численно заданного исходного распределения вероятностей гипотез;
  \item доказательной силы конкретных наблюдений, рассчитываемой независимо от ответа LLM;
  \item одинаковых наборов данных во всех сравниваемых условиях;
  \item последующего сравнения сырых наблюдений и статистического свидетельства, полученного из тех же наблюдений.
\end{enumerate}

Такой дизайн позволяет перейти от вопроса «может ли модель изменить ответ» к вопросу «соответствует ли вероятность правильного пересмотра количественной силе свидетельства и на каком этапе возникает расхождение».

\section{Методика первого эксперимента}

\subsection{Физическая система и конкурирующие гипотезы}

Для первого эксперимента выбрана простая синтетическая система с линейным законом и его нелинейным расширением. Нулевая гипотеза задаётся как

\begin{equation}
  \Hnull: \quad y = kx + \varepsilon,
\end{equation}

а альтернативная гипотеза -- как

\begin{equation}
  \Hone: \quad y = kx + \alpha x^3 + \varepsilon,
\end{equation}

где $k=1$, $\alpha$ определяет выраженность нелинейной компоненты, а шум $\varepsilon$ имеет нормальное распределение с нулевым математическим ожиданием. Гипотезы являются вложенными: при $\alpha=0$ модель \Hone{} совпадает с \Hnull. Поэтому сравнение только по ошибке подгонки некорректно: более сложная гипотеза может уменьшить ошибку даже на данных, порождённых \Hnull.

Для отрицательного контроля генерируются наблюдения при истинной \Hnull, а для исследования обнаружения нелинейности -- наблюдения при истинной \Hone. Это позволяет оценивать не только способность отказаться от ошибочной линейной гипотезы, но и склонность необоснованно выбирать более сложную модель.

\subsection{Объективная доказательная сила}

Основной мерой поддержки гипотез для конкретного набора наблюдений выбрана разность байесовских информационных критериев. Для модели $H$ критерий рассчитывается по формуле~\cite{schwarz}

\begin{equation}
  \mathrm{BIC}_H = n \ln\!\left(\frac{\mathrm{SSE}_H}{n}\right) + p_H \ln n,
\end{equation}

где $n$ -- число наблюдений, $\mathrm{SSE}_H$ -- сумма квадратов отклонений, а $p_H$ -- число оцениваемых параметров. В реализации используется

\begin{equation}
  \DeltaBIC = \mathrm{BIC}_{H_0} - \mathrm{BIC}_{H_1}.
\end{equation}

Положительное значение \DeltaBIC{} означает, что после учёта штрафа за дополнительный параметр данные лучше поддерживают \Hone; отрицательное значение поддерживает \Hnull. В качестве дополнительной диагностической характеристики сохраняется логарифм отношения правдоподобий, однако он не используется как единственный критерий выбора из-за вложенности моделей.

Следует различать доказательную силу конкретной выборки и калиброванный уровень режима генерации. Для комбинаций диапазона $x$, относительной нелинейности и шума методом Монте-Карло оценивается вероятность того, что правило на основе \DeltaBIC{} поддержит истинную гипотезу. По условию \Hone{} выбираются режимы, ближайшие к двенадцати целевым уровням от $0{,}50$ до $1{,}00$. Для тех же режимов отдельно рассчитывается вероятность корректной поддержки \Hnull, но она используется как диагностическая характеристика, а не как второй целевой параметр.

\subsection{Факторы и контроли}

Первый эксперимент имеет сбалансированный факторный дизайн. Его основные факторы представлены в таблице~\ref{tab:factors}.

\begin{table}[ht]
  \centering
  \caption{Факторы первого эксперимента}
  \label{tab:factors}
  \begin{tabularx}{\textwidth}{>{\raggedright\arraybackslash}p{0.26\textwidth} X}
    \toprule
    \textbf{Фактор} & \textbf{Уровни} \\
    \midrule
    Истинная гипотеза & \Hnull{} и \Hone{} \\
    Исходное убеждение & нейтральное: $P(H_0)=P(H_1)=0{,}5$; ошибочное: $0{,}9$ в пользу неверной гипотезы \\
    Доказательная сила & 12 калиброванных режимов генерации \\
    Порядок предъявления & \Hnull/\Hone{} и \Hone/\Hnull{} \\
    Языковая модель & Qwen3-0.6B и Qwen3-4B-AWQ \\
    \bottomrule
  \end{tabularx}
\end{table}

Для каждого режима подготовлено по 50 независимых наборов с истинной \Hnull{} и \Hone, всего 1200 наборов. Один набор используется в обоих условиях исходного убеждения и для обеих LLM. Порядок гипотез контрбалансируется внутри комбинаций режима и истинной гипотезы. Такая парная организация уменьшает влияние случайных различий между выборками и позволяет связывать изменение ответа именно с экспериментальным условием.

Для каждой модели предусмотрено 2400 испытаний, а общий объём для двух моделей составляет 4800 испытаний. Каждый запрос выполняется в независимом контексте. Температура генерации равна нулю, режим рассуждения отключён, а версии весов фиксируются точным идентификатором ревизии. Локальный запуск выполняется через vLLM, использующий механизм PagedAttention для эффективного обслуживания моделей~\cite{vllm}.

\subsection{Метрики и процедура анализа}

Основная зависимая переменная -- выбранная моделью гипотеза. Для каждого условия рассчитываются accuracy и balanced accuracy, доля выбора \Hone, true positive rate, true negative rate, false positive rate и false negative rate. Для долей рассчитываются 95\%-ные доверительные интервалы Уилсона~\cite{wilson}. Отдельные разрезы формируются по истинной гипотезе, исходному убеждению, порядку предъявления и уровню доказательной силы.

Ключевым объектом анализа является не одна усреднённая точность, а форма зависимости вероятности правильного выбора от доказательной силы. Сравнение нейтрального и ошибочного исходного убеждения на одних и тех же данных позволяет оценить сдвиг этой зависимости. Контроль с истинной \Hnull{} нужен для того, чтобы постоянный выбор более сложной гипотезы не интерпретировался как успешный пересмотр.

Первый эксперимент одновременно предъявляет модели сырые наблюдения и рассчитанные статистические характеристики. Он предназначен для ответа на основной исследовательский вопрос, однако вырожденный постоянный выбор \Hone{} не позволяет оценить форму зависимости от доказательной силы и требует предварительных абляций стимула. Для ответа на дополнительный вопрос потребуется отдельное сравнение условий \textit{raw data only} и \textit{processed evidence only} на тех же выборках.

\section{Основные результаты практики}

\subsection{Теоретические и методологические результаты}

В ходе практики проанализированы восемь близких работ по трём направлениям: пересмотр убеждений и контекстные смещения, оценка свидетельств языковыми моделями, а также использование LLM в интерактивном научном поиске. Для каждой работы выделены постановка, метод, основные выводы, ограничения и связь с настоящим исследованием. По итогам анализа сформулирован узкий исследовательский пробел, не сводящийся к общему утверждению о недостаточной надёжности LLM.

Зафиксированы основной и дополнительный исследовательские вопросы, операциональные определения исходного убеждения, доказательной силы и пересмотра, а также контрольные условия. Разработан подробный протокол первого эксперимента. В протоколе заранее определены факторы, объём выборки, правила контрбалансировки, основная статистическая мера и состав сохраняемых артефактов. Это уменьшает риск объяснять дизайн задним числом после просмотра ответов моделей.

\subsection{Реализованный экспериментальный конвейер}

В \href{https://github.com/drema3-4/Master-thesis.git}{репозитории проекта на GitHub} (см. также приложение~\ref{app:repository}) реализован сквозной программный конвейер первого эксперимента. Его основные этапы:

\begin{enumerate}
  \item генерация синтетических наблюдений для \Hnull{} и \Hone{} на фиксированной сетке параметров;
  \item Монте-Карло калибровка режимов по вероятности корректного знака \DeltaBIC;
  \item выбор режимов, соответствующих целевым уровням доказательной силы, и создание независимых наборов с уникальными начальными состояниями генератора;
  \item формирование запроса с контрбалансированным порядком гипотез и численно заданным исходным распределением вероятностей;
  \item обращение к локальной модели через OpenAI-совместимый интерфейс vLLM;
  \item строгий разбор ответа, фиксация ошибок и сохранение результата каждого испытания;
  \item автоматическое формирование сводных метрик и разрезов по факторам эксперимента.
\end{enumerate}

Логика разделена на модули генерации данных, вычисления статистических характеристик, сборки испытаний, взаимодействия с LLM, схем данных и служебных отчётов. Параметры вынесены в YAML-конфигурации и проверяются типизированной схемой \code{Ex1Config}. Она объединяет описание гипотез, параметры калибровки и набора данных, факторы эксперимента, настройки модели и генерации, пути к запросам и каталог результатов. Полная \hyperref[fig:config-schema]{графическая схема конфигурации} приведена в приложении~\ref{app:config-schema}.

\subsection{Воспроизводимость и контроль качества}

Для каждого запуска предусмотрена отдельная директория со следующими артефактами:

\begin{itemize}
  \item исходные ответы модели и параметры соответствующих испытаний;
  \item журнал ошибок, пустых или неразобранных ответов;
  \item разрешённая конфигурация запуска;
  \item манифест с идентификатором модели и входными файлами;
  \item сводные метрики и разрезы по экспериментальным факторам.
\end{itemize}

Калибровочный и основной наборы используют разные начальные состояния генератора случайных чисел. Для каждой конфигурации фиксируются путь к запросам, параметры генерации, версия схемы данных, идентификатор модели, ревизия весов и Git-коммит. Наборы \Hnull{} и \Hone{} сбалансированы, а их порядок в запросах контрбалансирован. Такие решения позволяют отличать содержательный эффект от повторения данных, позиционного смещения и изменений программного окружения. Связь входных артефактов, проверок и файлов результата показана на \hyperref[fig:repro-schema]{схеме воспроизводимости и контроля качества} в приложении~\ref{app:repro-schema}.

Подготовлены две конфигурации локальных моделей: Qwen3-0.6B как компактная базовая модель и Qwen3-4B-AWQ как более мощная квантованная модель. Для них используются одинаковые данные и параметры формирования ответа. Конфигурация контейнерного запуска учитывает доступные вычислительные ресурсы, что делает эксперимент воспроизводимым на используемой рабочей станции без обращения к внешнему закрытому API.

\subsection{Результаты запуска и анализа первого эксперимента}

\paragraph{Техническая полнота запусков}

Анализ выполнен в ноутбуке \path{notebooks/ex1_v4_short_analysis.ipynb} по последним завершённым запускам версии v4:

\begin{itemize}
  \item Qwen3-0.6B -- \path{ex1_qwen3_06b_v4_20261007T115433Z_5b4db7f8};
  \item Qwen3-4B-AWQ -- \path{ex1_qwen3_4b_awq_v4_20261007T121225Z_1c1239e7}.
\end{itemize}

Каждый запуск содержит 2400 валидных ответов, ноль ошибок выполнения и ноль неразобранных ответов. Продолжительность составила 9,4 минуты для Qwen3-0.6B и 38,6 минуты для Qwen3-4B-AWQ. Незавершённый каталог первой попытки запуска 4B-AWQ не включался в анализ. Числовые поля, видимые модели, совпадают с версией v3; в v4 изменена системная инструкция, явно объясняющая знак \DeltaBIC{}, вложенность гипотез и недопустимость выбора \Hone{} только по меньшей SSE или положительному отношению правдоподобий.

\paragraph{Основной количественный результат}

Обе модели во всех условиях выбрали \Hone{}: по 2400 ответов из 2400. Сводные значения приведены в таблице~\ref{tab:main-results}. Поскольку набор сбалансирован по истинным \Hnull{} и \Hone, постоянный ответ автоматически даёт 100\% точности на данных \Hone, 0\% на данных \Hnull{} и 50\% общей и сбалансированной точности. Следовательно, значение 50\% не следует интерпретировать как случайное угадывание: оно получено двумя одинаковыми константными стратегиями.

\begin{table}[ht]
  \centering
  \small
  \caption{Результаты последних завершённых запусков версии v4}
  \label{tab:main-results}
  \begin{tabularx}{\textwidth}{>{\raggedright\arraybackslash}X >{\raggedright\arraybackslash}p{0.14\textwidth} r r r r}
    \toprule
    \textbf{Модель} & \textbf{Prior} & \textbf{$n$} & \textbf{\Hone, \%} & \textbf{Accuracy, \%} & \textbf{Balanced acc., \%} \\
    \midrule
    Qwen3-0.6B & neutral & 1200 & 100,0 & 50,0 & 50,0 \\
    Qwen3-0.6B & wrong   & 1200 & 100,0 & 50,0 & 50,0 \\
    Qwen3-4B-AWQ & neutral & 1200 & 100,0 & 50,0 & 50,0 \\
    Qwen3-4B-AWQ & wrong   & 1200 & 100,0 & 50,0 & 50,0 \\
    \bottomrule
  \end{tabularx}
\end{table}

\paragraph{Различимость данных и конфликт статистических признаков}

Вырожденный ответ нельзя объяснить тем, что гипотезы неразличимы по данным. При нейтральном исходном распределении знак \DeltaBIC{} поддерживает истинную гипотезу в 1032 из 1200 уникальных наборов, то есть в 86\% случаев. В 722 наборах \DeltaBIC{} отрицателен и BIC поддерживает \Hnull, но модели всё равно выбирают \Hone. Диагностические характеристики приведены в таблице~\ref{tab:diagnostics}.

\begin{table}[ht]
  \centering
  \small
  \caption{Диагностика различимости гипотез на 1200 уникальных наборах}
  \label{tab:diagnostics}
  \begin{tabularx}{\textwidth}{X r r}
    \toprule
    \textbf{Условие} & \textbf{Число наборов} & \textbf{Доля, \%} \\
    \midrule
    Знак \DeltaBIC{} поддерживает истинную гипотезу & 1032 & 86,00 \\
    $\Delta\mathrm{BIC}<0$, BIC поддерживает \Hnull{} & 722 & 60,17 \\
    $\mathrm{SSE}_{H_1}<\mathrm{SSE}_{H_0}$ & 1200 & 100,00 \\
    $\log (L_{H_1}/L_{H_0})>0$ & 1200 & 100,00 \\
    \bottomrule
  \end{tabularx}
\end{table}

Последние две строки объясняются вложенностью моделей: более гибкая \Hone{} после подгонки на тех же данных всегда уменьшает SSE и увеличивает правдоподобие, тогда как BIC дополнительно штрафует её за параметр $\alpha$. Усиленная инструкция версии v4 предупреждает об этом конфликте, но не устраняет постоянный выбор \Hone. Наиболее правдоподобная интерпретация состоит в том, что модель использует простой однонаправленный признак -- положительное отношение правдоподобий, расположенное последней строкой запроса, либо меньшую ошибку более сложной модели -- и игнорирует штраф BIC. Это гипотеза о механизме, а не доказанный причинный вывод.

\paragraph{Влияние исходного убеждения, силы свидетельства и порядка}

Для каждого набора сравнивались нейтральное и намеренно ошибочное исходные распределения вероятностей. Нормативный выбор на основе prior и $\Delta\mathrm{BIC}/2$ должен измениться в 724 из 1200 пар, тогда как ни одна модель не изменила ответ ни в одной паре. Поэтому совпадение константного ответа с ожидаемым выбором, равное 39,83\% для neutral и 75,67\% для wrong, не является свидетельством байесовской чувствительности: различие возникает из-за распределения ожидаемых меток.

Для ветки с истинной \Hone{} нормативная средняя апостериорная вероятность в целом растёт с целевой доказательной силой и зависит от исходного распределения. Фактическая доля выбора \Hone{} остаётся равной 100\% на всех двенадцати уровнях. Контрбалансировка порядка также не меняет ответ: \Hone{} выбирается в 100\% случаев и в порядке \Hnull/\Hone, и в порядке \Hone/\Hnull. Это исключает простое позиционное смещение, но не отделяет склонность к метке \Hone{} от склонности к более гибкой формуле.

\paragraph{Воспроизводимость наблюдаемого паттерна}

Результат устойчив между повторными запусками, размерами моделей и версиями запроса (таблица~\ref{tab:repro-results}). Почти полное совпадение с v3 при существенно усиленной инструкции v4 показывает, что дальнейшее добавление словесных предупреждений без изменения структуры стимула вряд ли устранит ошибку.

\begin{table}[ht]
  \centering
  \small
  \caption{Согласованность ответов между запусками}
  \label{tab:repro-results}
  \begin{tabularx}{\textwidth}{X r r r}
    \toprule
    \textbf{Сравнение} & \textbf{Совпало} & \textbf{Из} & \textbf{Доля, \%} \\
    \midrule
    Два повтора Qwen3-0.6B, v4 & 2399 & 2400 & 99,96 \\
    Последние Qwen3-0.6B и Qwen3-4B-AWQ & 2400 & 2400 & 100,00 \\
    Qwen3-0.6B: v3 и последний v4 & 2397 & 2400 & 99,88 \\
    Qwen3-4B-AWQ: v3 и последний v4 & 2399 & 2400 & 99,96 \\
    \bottomrule
  \end{tabularx}
\end{table}

Таким образом, запуски технически корректны, но текущая постановка содержательно не позволяет оценить градиент пересмотра: обе модели используют одну и ту же неадекватную стратегию. Увеличение модели примерно в четыре раза повысило время запуска, но не изменило поведение. Полученный отрицательный результат локализует проблему в представлении конфликтующих статистических признаков и задаёт конкретные абляции для следующей версии эксперимента.

\section{Направления дальнейшей работы}

Ближайший этап должен проверить причину постоянного выбора \Hone{} с помощью минимальных абляций. Сначала модели следует предъявить только prior, \DeltaBIC{} и явно вычисленное логарифмическое апостериорное отношение

\begin{equation}
  \log\frac{P(H_1\mid D)}{P(H_0\mid D)}
  = \log\frac{P(H_1)}{P(H_0)} + \frac{\Delta\mathrm{BIC}}{2},
\end{equation}

расположив итоговую величину последней строкой запроса. Затем необходимо по одному возвращать SSE, значения правдоподобия и их отношение. Такое сравнение напрямую проверит гипотезу о конфликтующем либо последнем числовом признаке. Отдельный label-swap с обозначениями A/B позволит отделить склонность к метке \Hone{} от предпочтения более гибкой формулы.

После устранения вырожденного ответа следует повторить сравнение нейтрального и ошибочного исходного убеждения, проверить режимы с включённым и отключённым рассуждением и только затем сравнивать более крупные модели. Любое изменение запроса, состава признаков или режима генерации должно оформляться как новая версия эксперимента с отдельной конфигурацией и каталогом запуска.

Следующий содержательный этап -- разделение форм свидетельства. В условии с сырыми данными модель должна самостоятельно извлечь статистическую информацию из наблюдений, а в условии с обработанным свидетельством получит рассчитанные характеристики без исходного массива. Сравнение на одних и тех же выборках позволит локализовать ошибку: на этапе обработки чисел или на этапе объединения явно заданного свидетельства с исходным убеждением. Дальнейшее расширение возможно за счёт второй физической системы, моделей других семейств, условия верного исходного убеждения и контроля с нерелевантными числами.

\section{Заключение}

В ходе практики сформирована теоретическая и программная основа исследования пересмотра физических гипотез большими языковыми моделями. Анализ литературы показал, что модели уже используются в задачах научного поиска, однако комплексные бенчмарки обычно не отделяют пересмотр гипотезы от выбора эксперимента и извлечения закономерности. Работы по belief revision и байесовской интерпретации уверенности подробно исследуют влияние нового контекста, но редко используют физические числовые наблюдения с независимо рассчитанной доказательной силой.

На основании выявленного пробела предложена контролируемая постановка с двумя вложенными гипотезами и численной шкалой свидетельства. Разработанный дизайн включает нейтральное и ошибочное исходные убеждения, данные с истинной \Hnull{} и \Hone, контрбалансировку порядка гипотез и повторное использование одинаковых наборов во всех сравниваемых условиях. Основной статистический ориентир определяется через \DeltaBIC, а режимы генерации калибруются методом Монте-Карло.

Практическим результатом стала реализация воспроизводимого экспериментального конвейера: генерации и калибровки данных, сборки запросов, локального запуска двух LLM, разбора ответов и сохранения полного набора артефактов. В последних завершённых запусках Qwen3-0.6B и Qwen3-4B-AWQ получено по 2400 валидных ответов без ошибок. Обе модели во всех испытаниях выбрали \Hone, поэтому 50\%-ная точность на сбалансированном наборе является следствием константной стратегии.

При этом данные содержат различимый нормативный сигнал: знак \DeltaBIC{} соответствует истинной гипотезе в 86\% уникальных наборов, а изменение prior должно менять решение в 724 из 1200 пар. Фактически ответы не меняются ни в одной паре и не зависят от уровня доказательной силы или порядка предъявления. Результат воспроизводится между моделями, повторами и версиями запроса. Наиболее вероятное объяснение связано с предпочтением однонаправленного признака -- меньшей SSE либо положительного отношения правдоподобий для вложенной более сложной модели, -- однако для причинного разделения этого механизма, смещения к метке и эффекта последней строки нужны абляции.

Поставленные задачи практики выполнены: сформирована теоретическая постановка, реализован и проверен экспериментальный конвейер, проведены основные запуски и выявлено ограничение текущего стимула. Следующим этапом является упрощение набора предъявляемых статистических признаков, label-swap и повторный запуск, после чего можно переходить к сравнению сырых данных и явно предъявленного статистического свидетельства.

\clearpage
\renewcommand{\refname}{Список использованных источников}
\addcontentsline{toc}{section}{Список использованных источников}
\begin{thebibliography}{99}

\bibitem{discoverphysics}
Wiemann M. L., Smith L. M., Melchior P., Mishra-Sharma S., Wilson A. G., Izmailov P., Cuesta-Lázaro C. DiscoverPhysics: Benchmarking LLMs for Out-of-the-Box Scientific Thinking~// arXiv preprint arXiv:2605.26087. -- 2026.

\bibitem{physgym}
Chen Y., Pi\k{e}kos P., Ostaszewski M., Laakom F., Schmidhuber J. PhysGym: Benchmarking LLMs in Interactive Physics Discovery with Controlled Priors~// Advances in Neural Information Processing Systems. -- 2025.

\bibitem{beliefrevision}
Wilie B., Cahyawijaya S., Ishii E., He J., Fung P. Belief Revision: The Adaptability of Large Language Models Reasoning~// Proceedings of the 2024 Conference on Empirical Methods in Natural Language Processing. -- 2024. -- P.~10480--10496.

\bibitem{competingbiases}
Kumaran D., Fleming S. M., Markeeva L. et al. Competing Biases underlie Overconfidence and Underconfidence in LLMs~// Nature Machine Intelligence. -- 2026. -- Vol.~8. -- P.~614--627. -- DOI: \href{https://doi.org/10.1038/s42256-026-01217-9}{10.1038/s42256-026-01217-9}.

\bibitem{anchoring}
Owusu H. N., Feldman N. H. Anchoring Depends on Confidence and Post-Training in Language Models~// Proceedings of the 64th Annual Meeting of the Association for Computational Linguistics. Vol.~2: Short Papers. -- 2026. -- P.~174--180.

\bibitem{evidencetobelief}
Kim M., Kim S., Thorne J. From Evidence to Belief: A Bayesian Epistemology Approach to Language Models~// Proceedings of the 2025 Conference of the Nations of the Americas Chapter of the Association for Computational Linguistics: Human Language Technologies. Vol.~1. -- 2025. -- P.~10578--10611.

\bibitem{experimentsoutcomes}
Mohammadi S., Gaur M., Ferraro F. Experiments or Outcomes? Probing Scientific Feasibility in Large Language Models~// Proceedings of the 64th Annual Meeting of the Association for Computational Linguistics. Vol.~2: Short Papers. -- 2026. -- P.~605--614.

\bibitem{aihilbert}
Cory-Wright R., Cornelio C., Dash S., El Khadir B., Horesh L. Evolving Scientific Discovery by Unifying Data and Background Knowledge with AI Hilbert~// arXiv preprint arXiv:2308.09474. -- 2024.

\bibitem{schwarz}
Schwarz G. Estimating the Dimension of a Model~// The Annals of Statistics. -- 1978. -- Vol.~6, No.~2. -- P.~461--464. -- DOI: \href{https://doi.org/10.1214/aos/1176344136}{10.1214/aos/1176344136}.

\bibitem{wilson}
Wilson E. B. Probable Inference, the Law of Succession, and Statistical Inference~// Journal of the American Statistical Association. -- 1927. -- Vol.~22, No.~158. -- P.~209--212. -- DOI: \href{https://doi.org/10.1080/01621459.1927.10502953}{10.1080/01621459.1927.10502953}.

\bibitem{vllm}
Kwon W., Li Z., Zhuang S. et al. Efficient Memory Management for Large Language Model Serving with PagedAttention~// Proceedings of the 29th Symposium on Operating Systems Principles. -- 2023. -- P.~611--626. -- DOI: \href{https://doi.org/10.1145/3600006.3613165}{10.1145/3600006.3613165}.

\end{thebibliography}

\clearpage
\appendix

\section{Ссылка на код проекта}
\label{app:repository}

Исходный код, конфигурации, запросы и описание процедуры запуска доступны в репозитории:

\begin{center}
  \url{https://github.com/drema3-4/Master-thesis.git}
\end{center}

\clearpage
\section{Схема воспроизводимости и контроля качества}
\label{app:repro-schema}

На рисунке~\ref{fig:repro-schema} показано, как версионированные входы преобразуются в проверяемый каталог запуска. Сплошные стрелки обозначают основной поток данных, пунктирные -- контрольные проверки соответствующего этапа.

\begin{figure}[ht]
  \centering
  \begin{tikzpicture}[
    flow/.style={draw, rounded corners=2pt, align=left, text width=8.4cm, minimum height=2.45cm, inner sep=6pt, fill=gray!6, font=\footnotesize},
    check/.style={draw, rounded corners=2pt, align=left, text width=4.7cm, minimum height=2.45cm, inner sep=6pt, fill=gray!16, font=\footnotesize},
    arrow/.style={-{Latex[length=2.5mm]}, line width=0.7pt},
    checkarrow/.style={-{Latex[length=2.2mm]}, dashed, line width=0.6pt}
  ]
    \node[flow] (inputs) at (-2.8,0) {\textbf{Версионированные входы}\\YAML-конфигурация, prompts, набор данных, seeds, Git-коммит, идентификатор и ревизия модели};
    \node[check] (inputcheck) at (4.65,0) {\textbf{Проверка входов}\\валидация \code{Ex1Config}, существование путей, фиксированные версии};

    \node[flow] (trials) at (-2.8,-2.85) {\textbf{Сборка испытаний}\\парное использование наборов, neutral/wrong prior, два порядка гипотез, ожидаемый выбор};
    \node[check] (trialcheck) at (4.65,-2.85) {\textbf{Проверка дизайна}\\1200 уникальных наборов, баланс \Hnull/\Hone, контрбалансировка порядка};

    \node[flow] (runtime) at (-2.8,-5.70) {\textbf{Локальный запуск vLLM}\\отдельный контекст, температура 0, фиксированный seed, режим рассуждения отключён};
    \node[check] (runtimecheck) at (4.65,-5.70) {\textbf{Проверка выполнения}\\2400 запланированных испытаний на модель, контроль завершения};

    \node[flow] (parser) at (-2.8,-8.55) {\textbf{Строгий разбор ответа}\\допустимые ответы \Hnull/\Hone, фиксация пустых, ошибочных и неразобранных результатов};
    \node[check] (parsercheck) at (4.65,-8.55) {\textbf{Проверка качества}\\2400 успешных ответов, 0 failures, 0 invalid в каждом основном запуске};

    \node[flow] (artifacts) at (-2.8,-11.40) {\textbf{Каталог запуска}\\\code{responses.jsonl}, \code{failures.jsonl}, \code{resolver\_config.yaml}, \code{manifest.json}, \code{metrics.json}};
    \node[check] (artifactcheck) at (4.65,-11.40) {\textbf{Проверка происхождения}\\run ID, timestamps, Git-коммит, модель, ревизия, входные файлы и число ответов};

    \node[flow] (analysis) at (-2.8,-14.25) {\textbf{Ноутбук анализа}\\агрегации по prior, истинной гипотезе, силе свидетельства и порядку; сравнение повторов};
    \node[check] (analysischeck) at (4.65,-14.25) {\textbf{Проверка устойчивости}\\99,88--100\% согласованности между версиями, повторами и моделями};

    \draw[arrow] (inputs) -- (trials);
    \draw[arrow] (trials) -- (runtime);
    \draw[arrow] (runtime) -- (parser);
    \draw[arrow] (parser) -- (artifacts);
    \draw[arrow] (artifacts) -- (analysis);
    \draw[checkarrow] (inputs) -- (inputcheck);
    \draw[checkarrow] (trials) -- (trialcheck);
    \draw[checkarrow] (runtime) -- (runtimecheck);
    \draw[checkarrow] (parser) -- (parsercheck);
    \draw[checkarrow] (artifacts) -- (artifactcheck);
    \draw[checkarrow] (analysis) -- (analysischeck);
  \end{tikzpicture}
  \caption{Воспроизводимость и контроль качества экспериментального запуска}
  \label{fig:repro-schema}
\end{figure}

\clearpage
\section{Схема конфигурации эксперимента}
\label{app:config-schema}

Корневая схема \code{Ex1Config}, реализованная средствами Pydantic, проверяет YAML-конфигурацию до начала генерации данных или обращения к модели. Её основные группы полей показаны на рисунке~\ref{fig:config-schema}; фактические значения каждого запуска дополнительно сохраняются в \code{resolver\_config.yaml}.

\begin{figure}[ht]
  \centering
  \begin{tikzpicture}[
    root/.style={draw, rounded corners=2pt, align=center, text width=6.0cm, minimum height=1.3cm, inner sep=6pt, fill=gray!22, font=\bfseries},
    group/.style={draw, rounded corners=2pt, align=left, text width=6.6cm, minimum height=4.6cm, inner sep=7pt, fill=gray!7, font=\footnotesize},
    arrow/.style={-{Latex[length=2.5mm]}, line width=0.7pt}
  ]
    \node[root] (root) {\code{Ex1Config}\\корневая схема YAML-конфигурации};

    \node[group] (science) at (-3.75,-3.5) {\textbf{Постановка и данные}\\[0.35em]
      \code{experiment\_id}, \code{schema\_version}\\
      \code{H0}, \code{H1}\\
      \code{dataset\_calibration\_params}\\
      \code{target\_evidence\_strength}\\
      \code{dataset\_seed}\\
      \code{n\_datasets\_per\_level}\\
      \code{true\_hypotheses}\\
      \code{hypothesis\_orders}};

    \node[group] (model) at (3.75,-3.5) {\textbf{Модель и генерация}\\[0.35em]
      \code{llm\_config}: repo ID, revision, local dir, served name, runtime image\\[0.35em]
      \code{client\_config}: base URL, API key, model\\[0.35em]
      \code{generation\_config}: temperature, top-p, top-k, max tokens, seed, thinking};

    \node[group] (prompts) at (-3.75,-9.0) {\textbf{Prompts и prior}\\[0.35em]
      \code{system\_prompt\_path}\\
      \code{prior\_prompts\_paths}\\
      \code{prior\_probabilities}\\[0.6em]
      Проверяются имена условий, пути к шаблонам и распределения вероятностей для истинных \Hnull{} и \Hone.};

    \node[group] (paths) at (3.75,-9.0) {\textbf{Пути и сохранение результата}\\[0.35em]
      \code{calibration\_dataset\_path}\\
      \code{choose\_dataset\_params\_path}\\
      \code{dataset\_path}\\
      \code{save\_run\_experiment\_path}\\[0.6em]
      Разрешённая конфигурация сохраняется рядом с ответами, манифестом и метриками.};

    \draw[arrow] (root) -- (science);
    \draw[arrow] (root) -- (model);
    \draw[arrow] (science) -- (prompts);
    \draw[arrow] (model) -- (paths);
  \end{tikzpicture}
  \caption{Структура конфигурационной схемы первого эксперимента}
  \label{fig:config-schema}
\end{figure}

\end{document}
